% Source: https://artofproblemsolving.com/community/c3h2387182_a_triangle_problem
% Stable identity: aops-topic-2387182-q01
The angles of the triangle are $\angle A=80^\circ$, $\angle B=60^\circ$, $\angle C=40^\circ$.

\textbf{Step 1: $AB=BD=DC$.} Let line $BP$ meet side $AC$ at $D$; as $P$ is inside the triangle, $P$ lies between $B$ and $D$. In triangle $DBC$, $\angle DBC=40^\circ=\angle DCB$, so $DB=DC$, and the exterior angle is $\angle BDA=40^\circ+40^\circ=80^\circ$. In triangle $ABD$, $\angle BAD=80^\circ=\angle BDA$, so $BA=BD$. Hence $AB=BD=DC$, and $\angle BDC=100^\circ$.

\textbf{Step 2: the circumcentre $O$ of triangle $ABP$.} We have $\angle APB=180^\circ-10^\circ-20^\circ=150^\circ$. This inscribed angle stands on the arc $AB$ not containing $P$, so that arc measures $300^\circ$ and the arc $AB$ through $P$ measures $60^\circ$. Hence the central angle $\angle AOB=60^\circ$, and since $OA=OB$, triangle $ABO$ is equilateral: $OA=OB=OP=AB$. Also $O$ lies on the opposite side of line $AB$ from $P$ (because $P$ is on the minor arc).

\begin{asy}
usepackage("amsmath");
size(8cm);
defaultpen(fontsize(10pt));
pair d(real a){return dir(a);}
pair B=(0,0), C=(10,0);
real ab=10*Sin(40)/Sin(80);
pair A=ab*dir(60);
pair P=extension(B,B+dir(40),A,A+dir(250));
pair D=extension(B,P,A,C);
pair O=ab*dir(120);
draw(A--B--C--cycle,linewidth(0.9));
draw(A--P,linewidth(0.8)); draw(B--D,linewidth(0.8));
draw(A--O--B,gray); draw(O--P,blue); draw(O--D,blue+dashed); draw(P--C,red+linewidth(0.8));
draw(arc(O,ab,300,360),gray+dotted);
dot(A); dot(B); dot(C); dot(P); dot(D); dot(O);
label("$A$",A,N); label("$B$",B,SW); label("$C$",C,SE); label("$P$",P,dir(-95)); label("$D$",D,NE); label("$O$",O,W);
\end{asy}

\textbf{Step 3: two isosceles triangles.} Since $\angle OBA=60^\circ$, $\angle ABP=20^\circ$ and $O$, $P$ are on opposite sides of $AB$, we get $\angle OBP=80^\circ$. In the isosceles triangle $OBP$ ($OB=OP$) this gives $\angle OPB=80^\circ$. In triangle $OBD$ we have $OB=AB=BD$ and $\angle OBD=\angle OBP=80^\circ$, so $\angle BDO=\dfrac{180^\circ-80^\circ}{2}=50^\circ$.

\textbf{Step 4: a parallelogram.} Turning from ray $BD$ towards ray $BA$ by $20^\circ$ and then by $60^\circ$ more reaches ray $BO$; as $80^\circ<180^\circ$, the points $A$ and $O$ lie on the same side of line $BD$. The rays $PO$ and $DA$ make the same angle $\angle OPB=80^\circ=\angle ADB$ with the rays $PB$ and $DB$, which point the same way along line $BD$, and lie on the same side of it. Hence $PO\parallel DA$, with the same direction. As $D$ lies between $A$ and $C$, ray $DA$ is opposite to ray $DC$, so $\overrightarrow{PO}$ has the direction of $\overrightarrow{CD}$; and $PO=AB=DC$. Therefore $\overrightarrow{PO}=\overrightarrow{CD}$, and $PODC$ is a parallelogram. In particular $PC\parallel OD$.

\textbf{Step 5: the angle.} The vertices $O$ and $C$ of the parallelogram lie on opposite sides of its diagonal $PD$, so the alternate angles at $P$ and $D$ are equal: $\angle DPC=\angle PDO=\angle BDO=50^\circ$. In triangle $PDC$, $\angle PDC=\angle BDC=100^\circ$, so $\angle PCD=180^\circ-100^\circ-50^\circ=30^\circ$. Finally, since $D$ lies on $CA$,
\[\angle PCB=\angle ACB-\angle PCA=40^\circ-30^\circ=\boxed{10^\circ}.\]
